% =============================================================================
%  appendix_alternative_information_sets.tex
%  Numbers synced to the re-run output (tab_dralacbn_information_set_robustness,
%  tab_var_stability_info_sets, paper_numbers.csv). Changes vs previous draft:
%   - p_hat 3.18% -> 3.20%;
%   - robustness peaks: monetary 0.039 -> 0.030, uncertainty 0.023 -> 0.013
%     (real-side 0.033 and lighter 0.021 unchanged); R^2 / OOS RMSE refreshed;
%   - stability shares refreshed (uncertainty 0.2026 -> 0.5196, etc.);
%   - "monetary largest response" reworded (it is now comparable to, not above,
%     the baseline); table notes moved into minipages.
% =============================================================================

\subsection{Alternative VAR information sets}
\label{app:alternative_information_sets}
This appendix examines whether the credit-risk response depends on the
information set used in the VAR. The baseline specification in the main text
uses a parsimonious real-side VAR. We compare it with three alternatives: a
lighter real-side specification, a monetary specification, and an uncertainty
specification; their composition is given in
Table~\ref{tab:variables_macro_all}. In all exercises, the Merton--Vasicek
credit-risk parameters are kept fixed at their baseline values
($\widehat p=3.20\%$, $\widehat\rho=0.051$). Differences across
specifications therefore come only from the macro-financial transmission
mechanism and from the selected satellite bridge from macro-financial
variables to the systematic credit factor.
Figure~\ref{fig:robust_var_irf_alternatives} reports the macro-financial
responses for the three alternative information sets. The lighter real-side
specification checks that the baseline response does not rely on the full
real-side block used in the main text. The monetary and uncertainty
specifications allow the GPR innovation to propagate through other dimensions
of the aggregate environment. Across specifications, the shock remains
persistent and the macro-financial responses are consistent with a
contractionary geopolitical-risk disturbance, although the strength and
timing of the propagation differ across information sets.
\begin{figure}[!htbp]
\centering
\begin{subfigure}{0.48\textwidth}
    \centering
    \includegraphics[width=\linewidth]{figures/robustness/real_side_light/var_irf_1sd.pdf}
    \caption{Lighter real-side information set}
    \label{fig:robust_var_irf_real_side_light}
\end{subfigure}
\hfill
\begin{subfigure}{0.48\textwidth}
    \centering
    \includegraphics[width=\linewidth]{figures/robustness/monetary/var_irf_1sd.pdf}
    \caption{Monetary information set}
    \label{fig:robust_var_irf_monetary}
\end{subfigure}
\vspace{0.4em}
\begin{subfigure}{0.48\textwidth}
    \centering
    \includegraphics[width=\linewidth]{figures/robustness/uncertainty/var_irf_1sd.pdf}
    \caption{Uncertainty information set}
    \label{fig:robust_var_irf_uncertainty}
\end{subfigure}
\caption{\textbf{Macro-financial responses under alternative VAR information
sets.}
\textit{Notes: Posterior-median generalized impulse responses to a
one-standard-deviation innovation in the GPR equation, with $68\%$ (dark) and
$90\%$ (light) pointwise credible bands. Each panel uses a different VAR
information set. The baseline real-side responses are reported in the main
text (Figure~\ref{fig:irfs_var}).}}
\label{fig:robust_var_irf_alternatives}
\end{figure}
The differences in macro-financial propagation translate into different
credit-risk responses. Figure~\ref{fig:robust_pd_girf_alternatives} reports
the portfolio default-probability responses for the three alternative
information sets, and Table~\ref{tab:dralacbn_information_set_robustness}
summarizes the peaks. All specifications generate a positive median response
of portfolio PDs to a GPR innovation. Relative to the baseline peak of
$+0.033$ percentage points at $h=3$ (Section~\ref{sec:empirical_results}),
the lighter real-side specification produces a smaller and earlier response
($+0.021$ percentage points at $h=1$), the monetary specification a comparable
one ($+0.030$ at $h=1$), and the uncertainty specification a smaller one
($+0.013$ at $h=1$). The positive credit-risk effect of geopolitical-risk
innovations is therefore not confined to the baseline real-side channel,
although its magnitude and timing depend on the macro-financial channels
included in the VAR.
\begin{figure}[!htbp]
\centering
\begin{subfigure}{0.31\textwidth}
    \centering
    \includegraphics[width=\linewidth]{figures/robustness/real_side_light/pd_girf_1sd.pdf}
    \caption{Lighter real-side}
    \label{fig:robust_pd_girf_real_side_light}
\end{subfigure}
\hfill
\begin{subfigure}{0.31\textwidth}
    \centering
    \includegraphics[width=\linewidth]{figures/robustness/monetary/pd_girf_1sd.pdf}
    \caption{Monetary}
    \label{fig:robust_pd_girf_monetary}
\end{subfigure}
\hfill
\begin{subfigure}{0.31\textwidth}
    \centering
    \includegraphics[width=\linewidth]{figures/robustness/uncertainty/pd_girf_1sd.pdf}
    \caption{Uncertainty}
    \label{fig:robust_pd_girf_uncertainty}
\end{subfigure}
\caption{\textbf{Portfolio default-probability responses under alternative VAR
information sets.}
\textit{Notes: Posterior-median generalized impulse responses of portfolio
default probabilities to a one-standard-deviation innovation in the GPR
equation, with $68\%$ (dark) and $90\%$ (light) pointwise credible bands, in
percentage points. The Merton--Vasicek parameters are fixed across
specifications; only the VAR information set and the selected satellite
bridge change. The baseline real-side response is reported in the main text
(Figure~\ref{fig:girf_Z_PD}).}}
\label{fig:robust_pd_girf_alternatives}
\end{figure}
\begin{table}[!htbp]
  \centering
  \scriptsize
  \tablestretch
  \compacttablecols
  \setlength{\tabcolsep}{6pt}
  \caption{\textbf{DRALACBN credit-risk responses across VAR information sets.}}
  \label{tab:dralacbn_information_set_robustness}
  \begin{minipage}{0.92\textwidth}
  \centering
  \resizebox{\textwidth}{!}{%
    \begin{tabular}{lcccc}
      \toprule
      Information set & $k$ & Sat. $R^2$ & OOS RMSE & Peak $\Delta$PD, one s.d. (pp) \\
      \midrule
      Real-side & 6 & 0.889 & 0.4203 & $0.033$ $(h=3)$ \\
      Real-side (lighter) & 4 & 0.865 & 0.4742 & $0.021$ $(h=1)$ \\
      Monetary & 5 & 0.633 & 0.8759 & $0.030$ $(h=1)$ \\
      Uncertainty & 5 & 0.769 & 0.6862 & $0.013$ $(h=1)$ \\
      \bottomrule
    \end{tabular}%
  }
  \medskip
  \footnotesize
  \justifying
  \noindent \textbf{Note:} The Merton--Vasicek systematic factor $Z$, asset
  correlation $\rho$, and TTC PD are fixed within the DRALACBN application.
  Across information sets, only the VAR transmission mechanism and the
  selected satellite bridge change. Peak $\Delta$PD is the largest
  posterior-median PD response to a one-standard-deviation innovation in the
  GPR equation, in percentage points, with the corresponding horizon in
  parentheses. $Z$ is reconstructed from the FRED delinquency proxy and kept
  in its original Merton--Vasicek scale.
  \end{minipage}
\end{table}
Table~\ref{tab:var_stability_info_sets} reports posterior stability
diagnostics for the BVAR draws. Stable posterior draws are obtained for all
retained information sets, but the stable share varies across specifications.
The monetary specification is the most stable, with a stable share of
$0.8147$, whereas the uncertainty and baseline real-side specifications have
the lowest stable shares, equal to $0.5196$ and $0.5868$; the lighter
real-side specification lies in between, at $0.6379$. Conditional on
stability, the largest-modulus eigenvalues remain close to but below one
across all specifications, with posterior medians between $0.9910$ and
$0.9961$ and 95th percentiles below one. The robustness results should
therefore be read as conditional on stable posterior draws.
\begin{table}[!htbp]
  \centering
  \scriptsize
  \tablestretch
  \compacttablecols
  \setlength{\tabcolsep}{5pt}
  \caption{\textbf{Posterior stability diagnostics for the BVAR information sets.}}
  \label{tab:var_stability_info_sets}
  \resizebox{1\textwidth}{!}{%
    \begin{tabular}{lccccccc}
      \toprule
      Spec & $k$ & $p$ & Draws & Stable share & Stable draws
      & $|\lambda|_{\max}$ median & $|\lambda|_{\max}$ p95 \\
      \midrule
      Real-side
      & 6 & 2 & 30\,000 & 0.5868 & 17\,605 & 0.9961 & 0.9996 \\
      Real-side (lighter)
      & 4 & 2 & 30\,000 & 0.6379 & 19\,137 & 0.9910 & 0.9991 \\
      Monetary
      & 5 & 2 & 30\,000 & 0.8147 & 24\,442 & 0.9944 & 0.9992 \\
      Uncertainty
      & 5 & 2 & 30\,000 & 0.5196 & 15\,588 & 0.9952 & 0.9995 \\
      \bottomrule
    \end{tabular}%
  }
  \par
  \medskip
  \begin{minipage}{\textwidth}
  \footnotesize
  \noindent \textbf{Note:} $k$ is the number of variables in the VAR,
  including the GPR index, and $p$ is the lag order. Stable share is the
  fraction of posterior draws whose companion-matrix spectral radius is
  strictly below one. The last two columns report the posterior median and
  95th percentile of the largest-modulus eigenvalue over stable draws.
  \end{minipage}
\end{table}
Overall, the robustness exercise supports the main empirical conclusion. The
estimated effect of a GPR innovation on aggregate U.S. credit risk is modest
in absolute terms, but it is positive across the retained information sets.
The amplitude varies because the VAR determines the macro-financial
propagation of the shock, while the satellite equation determines which
propagated variables are transmitted to the Merton--Vasicek credit block.
This sensitivity is informative: the closed-form response can be recomputed
under alternative macro-financial information sets while keeping the
credit-risk block fixed.
